\section{O que anda junto com o quê}
\label{sec:correlacao}

A Figura~\ref{fig:correlacao} resume, num único mapa, o quanto cada par
de indicadores anda junto. A medida usada (correlação de Spearman) vai
de 0, quando dois indicadores não têm relação, a 1, quando crescem
sempre juntos. A mortalidade ajustada foi retirada do mapa por ser
estatisticamente idêntica à mortalidade geral (correlação de 1,000
entre as duas; todas as demais relações coincidem casa a casa).

Os pares mais fortes fazem sentido imediato para a gestão:
\textbf{custo por saída e fração de alta complexidade (0,63)}, custo e
tempo de permanência (0,56), e as duas ocupações entre si (0,48). A
\textbf{complexidade estrutural conversa com tudo}: 0,55 com alta
complexidade, 0,57 com ocupação de internação, 0,49 com custo, 0,47 com
TMP e 0,28 com mortalidade, o retrato numérico do papel de grande
divisor que a Seção~\ref{sec:complexidadeporte} descreveu. Já a
mortalidade se associa mais à ocupação de UTI (0,40) do que à de
internação (0,17), na direção esperada: onde a casuística é mais grave,
usa se mais UTI e morre se mais.

\begin{figure}[H]
\centering
\includegraphics[width=.96\textwidth]{fig_ae_10_correlacao_spearman.png}
\caption{Mapa de correlações entre os indicadores, as ocupações e a
complexidade estrutural. Azul mais escuro indica relação mais forte;
todos os valores estão impressos.}
\label{fig:correlacao}
\end{figure}

\comoler{procure a linha de um indicador e percorra as colunas: cada
célula diz o quanto aquele par anda junto (0 sem relação, 1 relação
perfeita). A diagonal vale sempre 1 porque compara o indicador com ele
mesmo. Repare na coluna da complexidade estrutural, escura quase
inteira: é o confundidor que precisa ser descontado em qualquer
comparação.}

A Figura~\ref{fig:variancia} e a Tabela~\ref{tab:variancia} respondem a
outra pergunta essencial: \textbf{quando um indicador varia, essa
variação vem de comparar hospitais diferentes ou do mesmo hospital
mudando ao longo dos anos?} A resposta é contundente: a variação vem
quase toda das diferenças entre hospitais. No tempo de permanência,
\textbf{apenas 2\% da variação total acontece dentro do mesmo hospital}
ao longo da década; na alta complexidade são 4,4\%, no custo 8,3\%, na
mortalidade 13,9\%, e só a ocupação de UTI passa de um terço (34,7\%).
Em termos de gestão: \textbf{cada hospital é muito parecido consigo
mesmo, ano após ano, e muito diferente dos vizinhos}. Os indicadores
têm forte inércia institucional, e mudanças reais dentro de um hospital
são eventos raros e valiosos de observar, o que dá ainda mais peso aos
cinco casos de conversão da Seção~\ref{sec:conversores}.

\begin{table}[H]
\centering
\caption{De onde vem a variação de cada indicador: parcela que ocorre
dentro do mesmo hospital ao longo dos anos e parcela que vem das
diferenças entre hospitais.}
\label{tab:variancia}
\small
\begin{tabular}{lrr}
\toprule
Indicador & Dentro do hospital (\%) & Entre hospitais (\%) \\
\midrule
TMP                    & 2,0  & 98,0 \\
Fração alta complex.   & 4,4  & 95,6 \\
Custo por saída        & 8,3  & 91,7 \\
Mortalidade geral      & 13,9 & 86,2 \\
Mortalidade ajustada   & 13,9 & 86,2 \\
Ocupação internação    & 17,9 & 82,1 \\
Ocupação UTI           & 34,7 & 65,3 \\
\bottomrule
\end{tabular}
\end{table}

\begin{figure}[H]
\centering
\includegraphics[width=.96\textwidth]{fig_ae_14_variancia_within.png}
\caption{Parcela da variação de cada indicador que acontece dentro do
mesmo hospital ao longo do tempo.}
\label{fig:variancia}
\end{figure}

\comoler{quanto menor a barra, mais estável o indicador dentro de cada
hospital e mais as diferenças observadas são traço permanente da
instituição. O TMP, com 2\%, é quase uma impressão digital do hospital;
a ocupação de UTI, com 35\%, é a que mais responde ao tempo, em boa
parte por causa da pandemia.}
